\cotefile{157-5}
\scanpage{405}
\mspage{319}
% archive sheet number: 404

% 1
On va distinguer \unsure{donc|dans} les deux cas.

% 2–4
2 A) $\alpha=\alpha_0+1$. Pour \unsure{construire} $i_\alpha$ et les $\tau_{\alpha,\beta}:i_\beta\unsure{-|\to}i_\alpha$, il suffit de construire $i_\alpha$ et \deleted{\unsure{les} $\tau$}
% 5
\[
\tikz[baseline=(c.base)]{\node[draw,ellipse,inner sep=2pt] (c) {$\tau_{\unsure{\alpha},\alpha_0}$};}: i_{\alpha_0}\longrightarrow i_\alpha=i_{\alpha_0+1}
\]

% 6–7
\unsure{D\,=}\ pour $\unsure{j|J}_\alpha$ et les $\tau'_{\deleted{\unsure{\beta}}\alpha,\unsure{\beta}}$, il suffit de construire
% 8
\[
\tikz[baseline=(c.base)]{\node[draw,ellipse,inner sep=2pt] (c) {$\tau'_{\alpha,\alpha_0}$};}:j_{\deleted{\unsure{\beta}}\alpha_0}\longrightarrow j_{\deleted{\unsure{\beta}}\alpha}.
\]

% 9
Il faut de plus se donner
% 10–11
\diagram{\begin{tikzcd}[ampersand replacement=\&,column sep=large,row sep=.3ex]
\deleted{\unsure{j}\;\unsure{\varphi(i_\alpha)}} \& {} \& {}\\[-1ex]
j_{\alpha_0}\arrow[r,"{v_{\alpha_0}}"{draw,ellipse,inner sep=1pt}] \& \varphi(\underset{\substack{\shortparallel\\i_{\alpha_0+1}}}{i_\alpha})\arrow[r,dash,"{u_\alpha}"{draw,ellipse,inner sep=1pt},"{\unsure{\,|\rangle}}"{at end,inner sep=0pt}] \& j_\alpha\;,
\end{tikzcd}}
% 12–14
donc en tout deux objets $i_\alpha,j_{\deleted{\unsure{\alpha}}\alpha}$ (dans $I$ et dans $\unsure{J'|I'}$ \unsure{respectivement}) et quatre flèches $\tau_{\alpha,\alpha_0},\tau'_{\alpha,\alpha_0},v_{\alpha_0},u_\alpha$.

% 15–20
Il faut quatre compatibilités\unsure{,|.} (page 318), où les deux dernières $3_\alpha(\beta)$, $\unsure{4}_\alpha(\beta)$ dépendent d'un paramètre $\beta<\alpha$ i.e. \uline{$\beta\leq\alpha_0$}, ou de $\beta<\alpha_0$.

% 21
\deleted{\unsure{Utilisant} que $\unsure{I}$}
\scanpage{406}
\mspage{320}
% archive sheet number: 405

\deleted{

% 1–4
Comme \deleted{$I$ est \unsure{filtrant|filtrante} et} $\varphi:I\unsure{-|\to}J$ est cofinal, \deleted{la catégorie $\unsure{j|J}_{\alpha_0}\backslash I$ est \unsure{pour} filtrante,} \inserted{il existe} un $i$ dans $I$ tel que $\varphi(i)$ \unsure{majore} $j_{\alpha_0}$, i.e. avec une flèche
% 5
\[
j_{\alpha_0}\longrightarrow\varphi(i).
\]

% 6–7
Comme $I$ est filtrante\unsure{,}, il existe un $i'$ dans $I$ et \unsure{des flèches}
% 8–9
\diagram{\begin{tikzcd}[ampersand replacement=\&]
i_{\alpha_0}\arrow[r] \& i'\,,\\
i\arrow[ur] \&
\end{tikzcd}}
% 10–11
donc quitte à remplacer $i$ par $i'$, \unsure{OPS|SPS} qu'il existe une flèche $\unsure{i|v}_{\alpha_0}\overset{\tau}{\longrightarrow}i$.

% 12–13
On \deleted{\unsure{prendra}} \inserted{\unsure{posera}} $i_\alpha=i$, et \deleted{$\unsure{i}_{\alpha_0}\longrightarrow i$} $\tau=\tau_{\unsure{\alpha},\alpha_0}$, \deleted{\unsure{enfin}} d'où
% 14
\begingroup\csname tagsleft@true\endcsname
\[
i_{\alpha_0}\overset{\tau_{\alpha,\alpha_0}}{\unsure{\text{---}|\longrightarrow}}i_\alpha
\tag*{\unsure{(a)|($\alpha$)|(*)}}
\]
\endgroup
% 15
et
% 16
\begingroup\csname tagsleft@true\endcsname
\[
j_{\alpha_0}\overset{v_{\alpha_0}}{\text{---}}\varphi(i_\alpha).
\tag*{\unsure{(aa)|(**)|(a')}}
\]
\endgroup

% 17–18
Comme $J'$ est cofinal \unsure{dans|sur} $J$, il existe $j_\alpha$ dans $J$ et une flèche
% 19
\begingroup\csname tagsleft@true\endcsname
\[
\varphi(i_\alpha)\xrightarrow{u_\alpha}j_{\unsure{\alpha|a}}.
\tag*{\unsure{(aaa)|(***)|(aa')}}
\]
\endgroup
}
\scanpage{407}
\mspage{32\deleted{\unsure{0}}1}
% archive sheet number: 406

% 1–2
\deleted{On va \unsure{maintenant} distinguer les deux cas

2 A) On a $\alpha=\alpha_0+1$.}

% 3–5
On va d'abord s'occuper des deux premières compatibilités $1_\alpha,2_\alpha$, en termes de \unsure{la} donnée de
% 6
\[
i_{\alpha_0}\in\operatorname{Ob} I,\quad \unsure{j|J}_{\alpha_0}\in\operatorname{Ob} J',\quad u_{\alpha_0}:\varphi(i_{\alpha_0})\longrightarrow j_{\alpha_0}.
\]

% 7–10
\deleted{\unsure{C|L}}\unsure{L|C}es \inserted{deux} compatibilités \unsure{voulues} pour les deux \unsure{objets} $i_\alpha,j_{\deleted{\unsure{\beta}}\alpha}$ et les quatre flèches $\tau_{\alpha,\alpha_0},\tau'_{\alpha,\alpha_0},v_{\alpha_0},u_\alpha$ sont indiquées dans le diagramme
% 11–13
\diagram{\begin{tikzcd}[ampersand replacement=\&,column sep=small]
\varphi(i_{\alpha_0})\arrow[r,"u_{\alpha_0}"]\arrow[rr,bend right=27,"{\varphi(\underline{\tau_{\alpha,\alpha_0}})}"'] \& j_{\alpha_0}\arrow[r,"{\underline{v_{\alpha_0}}\,?}"]\arrow[rr,bend left=27,"{\tau'_{\alpha,\alpha_0}}"] \& \varphi(\underset{?}{\underline{i_\alpha}})\arrow[r,"{\underline{u_\alpha}\,?}"] \& \underset{?}{j_\alpha}\qquad(*)
\end{tikzcd}}

% 14–16
On \unsure{note|voit} que comme $\varphi$ est cofinal, $\exists\,i_{\unsure{\alpha|?}}=i_{\deleted{\unsure{\alpha}}}\in\operatorname{Ob} I$ et $j_{\alpha_0}\xrightarrow{v_{\alpha_0}}\varphi(i_{\deleted{\unsure{\alpha}}})$. Mais on veut que le composé
% 17
\diagram{\begin{tikzcd}[ampersand replacement=\&,column sep=small]
w:\quad\varphi(i_{\alpha_0})\arrow[r,"u_{\alpha_0}"] \& j_{\alpha_0}\arrow[r,"v_{\alpha_0}"] \& \varphi(i_{\deleted{\unsure{\alpha}}})
\end{tikzcd}}
% 18–21
soit de la forme $\varphi(\tau)$, où $\tau=\tau_{\alpha,\alpha_0}:i_{\alpha_0}\longrightarrow i_\alpha$ est une flèche dans $I$. Comme $\varphi$ est un foncteur cofinal de catégories filtrantes, on sait
\scanpage{408}
\mspage{322}
% archive sheet number: 407

% 1–2
qu'on peut trouver $i\xrightarrow{\unsure{\rho|s}}i'$ dans $I$, tel que la flèche composée $\varphi(i_{\alpha_0})\xrightarrow{w}\varphi(i)\xrightarrow{\varphi(\rho)}\varphi(i')$
% 3–5
soit de la forme $\varphi(\tau)$. Donc quitte à remplacer $i$ par $i'$, et $v_{\alpha_0}$ par \unsure{la|le} \unsure{composé|composée} \inserted{$v'_{\alpha_0}:$} $\unsure{j|J}_{\deleted{\unsure{\beta}}\alpha_0}\xrightarrow{v_{\alpha_0}}\varphi(i)\xrightarrow{\varphi(\rho)}\varphi(i')$, \deleted{\unsure{OPS|SPS}}
% 6–10
\deleted{que} \unsure{OPS|SPS} que $v_{\alpha_0}u_{\alpha_0}$ est de la forme $\varphi(\tau)$. Cela nous donne donc en principe \unsure{l'}objet $i=i_\alpha$, et les flèches $v_{\alpha_0}:$ \inserted{$j_{\alpha_0}\to\varphi(i_{\unsure{\alpha|\tau}})$} et $\tau_{\alpha,\alpha_0}:i_{\alpha_0}\unsure{-|\to}i_\alpha$, de façon à satisfaire à la première compati\deleted{\unsure{ti}}bilité \unsure{dans|de} \unsure{(*)|(a)},

% 11–15
Il reste à construire $j_\alpha,\tau'_{\alpha,\alpha_0},u_\alpha$ de façon à satisfaire à la deuxième compati\deleted{\unsure{o}}bilité $(*)$. Or $J'$ étant cofinal \unsure{dans|sur} $J$, il existe $j_\alpha$ dans $J'$ et une flèche $u_\alpha:\varphi(i_\alpha)\longrightarrow\deleted{\unsure{J}}j_\alpha$.

% 16–20
D'autre part, comme $J'\unsure{-|\to}J$ est pl. fidèle, la flèche composée $j_{\alpha_0}\xrightarrow{v_{\alpha_0}}\varphi(i_\alpha)\xrightarrow{u_\alpha}j_\alpha$ dans $\unsure{\mathbb{J}|J|I}$, entre objets de $J'$, provient d'une unique flèche $\tau'_{\unsure{\alpha_0|\alpha},\alpha_0}:j_{\alpha_0}\longrightarrow j_\alpha$ dans $J'$.
\scanpage{409}
\mspage{323}
% archive sheet number: 408

% 1–6
Ainsi, on a construit les deux objets et les quatre flèches voulues, de façon à satisfaire aux \unsure{conditions|compatibilités|contraintes} $1_\alpha,2_{\unsure{\alpha}}$ de la page 318. Les contraintes $3_\alpha(\beta),4_\alpha(\beta)$ sont-elles alors \unsure{satisfaites} automatiquement ?
% 7–9
\diagram{\begin{tikzcd}[ampersand replacement=\&,column sep=large]
\varphi(i_\beta)\arrow[r,"u_\beta"]\arrow[d] \& \unsure{j|J}_\beta\arrow[d]\\
\varphi(i_{\alpha_0})\arrow[r,"u_{\alpha_0}"]\arrow[d,"{\varphi(\tau_{\alpha,\alpha_0})}"] \& \deleted{\unsure{\varphi(}}j_{\alpha_0}\deleted{\unsure{)}}\arrow[d,"{\tau'_{\alpha,\alpha_0}}"]\\
\varphi(i_\alpha)\arrow[r,"u_\alpha"] \& j_\alpha
\end{tikzcd}}

% 7–8, right
\margin{$\beta<\alpha$ i.e.\\ $\beta\leq\alpha_0$}

% 10–14
Le carré supérieur est commutatif par hyp. de récurrence, pour \deleted{\unsure{voir}} que le carré total \inserted{(composé)} soit commutatif, il suffit \deleted{de le \unsure{v}} que le carré inférieur le soit. On a donc une double-flèche
% 15
\[
\deleted{\unsure{\text{dans}}}\quad(\tau'_{\alpha,\alpha_0}u_{\alpha_0},u_\alpha\varphi(\tau_{\alpha,\alpha_0})):\varphi(i_{\alpha_0})\mathrel{\substack{\unsure{\hookrightarrow|\leftrightarrow|\rightarrow}\\[-.55ex]\rightarrow}}j_\alpha,
\]
% 16–19
elle n'est peut-être pas triviale mais on sait, $J'$ étant \unsure{filtrant|filtrante}, qu'on peut l'égaliser par une flèche $j_\alpha\to\overrightarrow{j}_\alpha$ dans $J'$. Donc quitte à remplacer \deleted{\unsure{$j$}}$j_\alpha$ par
\scanpage{410}
\mspage{324}
% archive sheet number: 409

% 1–5
\deleted{par} $\overline{\unsure{j|J}}_\alpha$, $u_\alpha$ par le composé $\overline{u}_\alpha:\varphi(i_{\unsure{\delta|\alpha}})\longrightarrow j_\alpha\to\overline{j}_\alpha$, et $\tau'_{\alpha,\alpha_0}$ par le composé $\overline{\tau'}_{\alpha,\alpha_0}:j_{\alpha_0}\to j_\alpha\to\overline{j}_{\unsure{\alpha_0|\alpha}}$, (sans toucher à $i_\alpha,v_{\alpha_0},\tau_{\alpha_0,\alpha}$), on peut obtenir la commutativité voulue, d'où \deleted{$\unsure{3_\alpha}$} $3_\alpha(\beta)$.

% 6–7
Reste la contrainte $4_\alpha(\beta)$. Pour ceci, considérons \unsure{le|les} \unsure{diagramme|diagrammes} \unsure{composite|composites}
% 8–10
\diagram{\begin{tikzcd}[ampersand replacement=\&,column sep=large]
j_\beta\arrow[r,"v_\beta"]\arrow[d,"{\tau'_{\alpha_0,\beta}}"'] \& \varphi(i_{\beta+1})\arrow[d,"{\varphi(\tau_{\alpha,\beta+1})}"]\\
j_{\alpha_0}\arrow[r,"v_{\alpha_0}"] \& \varphi(i_{\alpha_0\deleted{\unsure{\alpha}}+1})=\varphi(i_\alpha)\quad .
\end{tikzcd}}

% 11–17
\deleted{\unsure{Pour tout}} Si $\alpha_0$ n'est pas un ordinal limite, \deleted{mais} de la forme $\alpha_1+1$, on peut encore se ramener au cas où $\beta=\alpha_1$, en procédant comme ci-dessus. Mais on \deleted{peut} \unsure{procédera} dans tous les cas, \deleted{par une} sans un tel artifice (!), en considérant les doubles-flèches
% 18
\diagram{\begin{tikzcd}[ampersand replacement=\&,column sep=huge]
(\unsure{*})\qquad j_\beta\arrow[r,shift left=2.5pt,"{\varphi(\tau_{\alpha,\beta+1})v_\beta}"]\arrow[r,shift right=2.5pt,"{v_{\alpha_0}\tau'_{\alpha_0,\beta}}"'] \& \varphi(i_\alpha)
\end{tikzcd}}
% 19–20
déduites du carré précédent. Comme \deleted{$J'$ \unsure{est} $J$ est \unsure{grand} devant $\unsure{\pi|\Pi}$.}
\scanpage{411}
\mspage{\unsure{315|325}}
% archive sheet number: 410

% 1–6
\unsure{l'ens|l'ens.} de \unsure{ses|ces} doubles flèches est de cardinal $\operatorname{Card}\alpha<\unsure{\pi|\varpi}$, et que $\unsure{\mathbb{J}|J}$ est \unsure{grde|grande} devant $\pi$ (il suffirait qu'il le soit pour les cardinaux $\pi'<\pi$), on peut trouver une flèche $\varphi(i_\alpha)\xrightarrow{u}\unsure{j|J}$ dans $J$ qui égalise toutes ces doubles-flèches.

% 7–13
Comme $\varphi$ est cofinal, \unsure{OPS|SPS} que $j$ est de la forme $\varphi(i)$. Comme $\varphi$ est un foncteur cofinal de catégories filtrantes, il existe alors $i\overset{\unsure{\rho|p}}{\longrightarrow}i'$ dans $I$, tel que \deleted{$\unsure{\varphi(\rho)}$} le composé $\varphi(i_\alpha)\xrightarrow{u}\varphi(i)\xrightarrow{\varphi(\rho)}\varphi(i')$ soit de la forme $\varphi(\sigma)$, pour $\sigma:i_\alpha\longrightarrow i'$.
% 13–19
Au total, on a trouvé $\sigma:i_\alpha\longrightarrow i'$, tel que $\varphi(\sigma)\;(=\varphi(\rho)u):\varphi(i_\alpha)\unsure{-|\to}\varphi(i')$ égalise toutes les doubles-flèches envisagées dans $(\unsure{*})$ p.~324. Donc il s'impose de remplacer $i_\alpha$ par $\overline{i}_\alpha=i'$, $v_{\alpha_0}$ par le composé $j_{\alpha_0}\to\varphi(i_\alpha)\xrightarrow{\varphi(\sigma)}\varphi(\overline{i}_\alpha)$, $\tau_{\unsure{\alpha_0,\alpha|\alpha,\alpha_0}}$\unsure{.|,}
\scanpage{412}
\mspage{326}
% archive sheet number: 411

% 1
par \inserted{le composé} $\overline{\tau}_{\alpha,\alpha_0}:\deleted{\unsure{\alpha}}i_{\alpha_0}\xrightarrow{\tau_{\alpha,\alpha_0}}i_\alpha\xrightarrow{\sigma}\overline{i}_\alpha$.

% 2–4
Pour résumer, voici la recette pour construire, d'abord \tikz[baseline=(c.base)]{\node[draw,ellipse,inner sep=2pt] (c) {$i_\alpha,v_{\alpha_0},\tau_{\unsure{\alpha|\alpha_0},\alpha_0}$};}, et \uline{ensuite} \tikz[baseline=(c.base)]{\node[draw,ellipse,inner sep=2pt] (c) {$\unsure{j|J}_\alpha,u_\alpha,\tau'_{\alpha,\alpha_0}$};}.

% 5–9
On procède d'abord comme pages 321, 322 pour construire le premier triplet, de façon à satisfaire à la première compatibilité $(\unsure{*})$ (p.~321) i.e. à $1_\alpha$.
% 9–15
On ``corrige'' ensuite ce choix de $i_\alpha,v_{\alpha_0},\tau_{\alpha,\alpha_0}$, en \unsure{remplaçant} $i_\alpha$ par $\overline{i_\alpha}$ \unsure{par|pour} $\sigma:i_\alpha\to\overline{i}_\alpha$, $v_{\alpha_0}$ par $\overline{v}_{\alpha_0}=\varphi(\sigma)v_{\alpha_0}$, et $\tau_{\alpha,\alpha_0}$ par $\overline{\tau_{\alpha,\alpha_0}}=\sigma\unsure{-|\cdot|\,}\tau_{\alpha,\alpha_0}$, de sorte à obtenir les compatibilités $4_\alpha(\beta)$ pour $\beta<\alpha_0$, comme on vient de l'indiquer.
% 15–17
On ne touchera plus à ces choix de $i_\alpha,v_{\alpha_0},\tau_{\alpha,\alpha_0}$, dans la suite de la construction\unsure{.|,}

% 18–20
On construit $(j_\alpha,u_\alpha,\tau'_{\alpha,\alpha_0})$ de \unsure{façon} à satisfaire à la deuxième compatibilité $(*)$ (p.~321), comme
\scanpage{413}
\mspage{327}
% archive sheet number: 412

% 1–6
indiqué page 322. Ensuite, on corrige \unsure{ce|ces} choix, en remplaçant $\unsure{j|J}_\alpha$ par $\overline{j}_\alpha$, \deleted{\unsure{disons|alors}} \inserted{avec} $\rho:j_\alpha\to\overline{j}_\alpha$, et $u_\alpha$ par $\overline{u}_\alpha=\rho u_\alpha$, $\tau'_{\alpha,\alpha_0}$ par $\overline{\tau'}_{\alpha,\alpha_0}=\rho\tau'_{\alpha,\alpha_0}$, de façon à satisfaire aux contraintes \unsure{(}$3_\alpha$)$(\beta)$ de la page 318.

% 7–8
Cela achève la construction dans le cas où $\alpha=\alpha_0+1$. \deleted{\unsure{où}}

% 9
\uline{Cas 2 B} : \uline{$\alpha$ est un ordinal limite}.

% 10–13
Il faut donc construire les données 1) 2) 3) de la page 317 (à l'exclusion de la donnée 4), de façon à satisfaire à la seule contrainte $3_\alpha(\beta)$.

% 14–18
On commence à construire $i_\alpha$ et les $\tau_{\beta,\alpha}$ dans $I$. C'est possible, car $I_\alpha=\{\beta\mid<\beta\deleted{\unsure{\alpha}}<\alpha\}$ est de cardinal \deleted{$\unsure{\pi_\alpha}$} $\pi_\alpha<\unsure{\pi|\varpi}$ -- il suffit donc de supposer $I$ \unsure{grd|grand} devant les cardinaux \deleted{\unsure{petits}} $\pi'<\pi$,

% 19–20
On construit ensuite de \unsure{$=$} $j_\alpha$ \inserted{dans $J'$} \deleted{et} les $\tau'_{\beta,\alpha}$ $(\beta<\alpha)$, en utilisant \unsure{le|les}
\scanpage{414}
\mspage{\unsure{318|328}}
% archive sheet number: 413

% 1
fait que $J'$ est \unsure{grd|grand} devant \deleted{\unsure{Card}} $\operatorname{Card}\alpha=\operatorname{Card} I_\alpha$.

% 2–3
On doit de plus construire une flèche $u_\alpha:\varphi(i_\alpha)\longrightarrow\unsure{j|J}_\alpha$.

% 4–7
\deleted{\unsure{C}} C'est possible, car \deleted{il \unsure{exi}} $J'$ étant cofinale dans $J$, il existe $j'$ dans $J'$ et $\varphi(i_\alpha)\longrightarrow j'$, et $J'$ étant filtrante il existe $\overline{j_\alpha}$ et des flèches
% 8–9
\diagram{\begin{tikzcd}[ampersand replacement=\&]
j_\alpha\arrow[r] \& \overline{j_\alpha}\,;\\[-1ex]
j'\arrow[ur] \&
\end{tikzcd}}
% 9–12
donc quitte à remplacer $j_\alpha$ par $\overline{j}_\alpha$, les $\tau'_{\alpha,\beta}$ par des $\overline{\tau'_{\alpha,\beta}}$, on trouve bien une flèche \deleted{\unsure{$u_\alpha$}} $u_\alpha:\varphi(i_\alpha)\longrightarrow j_\alpha$. Il faut ensuite \deleted{vérifier les \unsure{cont}} s'occuper des compatibilités $3_\alpha(\beta)$, \deleted{Mais} donc des \deleted{\unsure{la}} doubles-flèches
% 13–15
\diagram{\begin{tikzcd}[ampersand replacement=\&,column sep=huge]
\varphi(i_\beta)\arrow[r,<-,shift left=2pt,"{\overline{\tau'}_{\alpha,\beta}u_\beta}","{\unsure{\,|\rangle}}"{at end,inner sep=0pt}]\arrow[r,<-,shift right=2pt,"{u_{\deleted{\unsure{\alpha}}\alpha}\varphi(\tau_{\alpha,\beta})}"'] \& j_\alpha\;.
\end{tikzcd}}

% 16–18
Comme \unsure{l'ens|l'ens.} de ces doubles flèches est de cardinal $\operatorname{Card}\alpha<\unsure{\pi|\varpi}$, et que $J^{\unsure{*|\prime}}$ est grd devant ledit cardinal, on
\scanpage{415}
\mspage{329}
% archive sheet number: 414

% 1
on trouve donc une flèche
% 2
\[
\unsure{\rho|p}:\unsure{j|J}_\alpha\longrightarrow\deleted{\unsure{\overline{j}_\alpha}}j
\]
% 3–5
dans $\unsure{\mathbb{J}|J}$ qui les égalise toutes. Comme $\deleted{\unsure{\varphi}}J'$ est cofinal \unsure{OPS|SPS} $j$ \deleted{de la forme $\varphi(i)$, \unsure{donc}} \inserted{dans $J'$, soit} $\overline{j}_\alpha$. Donc quitte
% 6–9
à remplacer $j_\alpha$ par $\overline{j}_\alpha$, et \deleted{les} les $\tau'_{\deleted{\unsure{\beta}}\alpha,\beta}$ par $\inserted{\overline{\tau}'_{\alpha,\beta}=}\rho\,\tau'_{\alpha\beta}$, enfin $u_\alpha:\varphi(i_\alpha)\to j_\alpha$ par $\overline{u}^{\unsure{\prime|\,}}_\alpha=\rho u_\alpha:\varphi(i_{\unsure{\alpha|b}})\to\overline{j}_\alpha$, on satisfait à toutes les contraintes voulues.

% 10–11
Je résume le résultat \inserted{(``$\pm$ tautologique'' !)} prouvé dans les pages 313--329 :

% 12–13
\llap{\rule[-1.0ex]{0.45pt}{3.1ex}\hspace{.7em}}\uline{Lemme}\inserted{-\unsure{déf|def}} 4.13.17. Considérons un carré cartésien dans $\mathrm{Cat}$
% 14–16
\diagram{\begin{tikzcd}[ampersand replacement=\&]
I\arrow[r,"\varphi"] \& J\\
I'\arrow[u,hook]\arrow[r] \& J'\arrow[u,hook,"\unsure{\psi|\varphi}"']
\end{tikzcd}}

% 17–19
On fait les \deleted{et un cardinal \inserted{\unsure{ordinal} \unsure{infini} limite $\alpha_0$ ; Soit $I_{\alpha_0}=\{\alpha\mid\alpha\ \text{ordinal}<\alpha_0\}$} $\pi$. On fait les} \llap{\rule[-1ex]{.45pt}{3ex}\hspace{.7em}}hypothèses suivantes :
\scanpage{416}
\mspage{330}
% archive sheet number: 415

\deleted{

% 1–2
2) \inserted{$\forall\,\alpha\in I_{\alpha_0}$,} $I$ est \unsure{grd|grand} devant \deleted{\unsure{tout cardinal} $\unsure{\pi'<\pi}$ \illegible{2}} \inserted{$\operatorname{Card}\alpha$, pour tt $\alpha\in I_{\alpha_0}$ ;} \deleted{(\unsure{ou} \illegible{3})}

% 3–4
1) \inserted{Le foncteur} $\varphi$ est cofinal. [NB On sait que \deleted{\unsure{cel}} \inserted{1) et 2)} \unsure{impliquent} que $J$ est grd devant \deleted{\unsure{tout cardinal} $\unsure{\pi'<\pi}$} \inserted{$\operatorname{Card}\alpha$, pour tt $\alpha\in I_{\alpha_0}$ ;} \inserted{(\unsure{resp.} \unsure{devant} $\unsure{\pi}$)} ]

% 5–8
2) $I$ est stable par $\varinjlim$ de type $I_{\alpha\deleted{\unsure{\pi}}}$ \deleted{(où $\alpha_\pi$ est le plus petit ordinal $\alpha$ \deleted{\unsure{tel}} que $\operatorname{card}\alpha=\pi$), i.e. par $\varinjlim$ indexées par des ens. \unsure{préord.} bien ordonnés} \deleted{\unsure{de cardinal} $\pi$}. et le foncteur $\varphi$ commute à ces limites.

% 9–13
3) Le foncteur $\psi$ est l'inclusion d'une sous-catégorie str. pleine $J'$ de $J$ ; \deleted{et $J'$ est} De plus $J'$ est stable par $\varinjlim$ du type \deleted{\unsure{envisagé}} $I_{\alpha_\pi}$, et \unsure{le|les} \unsure{foncteur|foncteurs} $\unsure{\psi|f}$ \unsure{commute|commutent} \unsure{aux} $\varinjlim$.

% 14–16
4) $J'$ est cofinale dans $J$. [Donc il est grand devant tt cardinal $\pi'<\pi$, \deleted{$J$ l'étant \unsure{(?)}} (resp. devant $\pi$) \unsure{$J$ l'étant.|$J$ l'étant aussi.}]

% 17–21
Sous ces conditions, $I'$ est cofinale dans $I$. [Et il s'ensuit aussitôt que $I'$, tout comme $I,J,J'$, est grd devant tt cardinal $\pi'<\pi$ (resp. devant $\pi$), et on voit aussitôt que le foncteur $I'\unsure{-|\to}J'$ est également cofinal.]
}
\scanpage{417}
\mspage{330}
% archive sheet number: 416

% 1–3
1) $\varphi$ et $\psi$ sont cofinaux, $\psi$ est l'inclusion d'une sous-catégorie str. pleine $J'$ de $J$, $I$ est filtrant (donc $J$ et $J'$ aussi).

% 4–10
On cherche des conditions pour que l'inclusion \deleted{\unsure{pl.}} de la sous-catégorie str. pleine $I'\unsure{-|\to}I$ soit cofinale (auquel cas on voit tout de suite que $I'$ est \unsure{filtrante|filtrant}, et $I'\to J'$ est aussi cofinal). Pour ceci, on suppose qu'il existe un \uline{ordinal limite} $\alpha_0$ ayant les propriétés \inserted{suppl.} suivantes\unsure{.|:}

% 11–13
2) $I$ et $J'$ sont stables par $\varinjlim$ du type $I_{\alpha_0}$ $(=\{\beta\mid\beta\ \text{ordinal}<\alpha_0\})$, et les foncteurs $\varphi,\psi$ commutent à ces limites.

% 14–18
3) Pour tt $\alpha\in I_{\alpha_0}$, $I$ est \deleted{\unsure{stab}} \unsure{grd|grand} devant $\operatorname{Card}\alpha$. [Donc, $\varphi$ étant cofinal, $J$ est \inserted{\unsure{égalemt|églmt|également}} grd devant $\operatorname{Card}\alpha$, et de $=$, $\psi$ étant pl. fidèle, et cofinal, $J'$ est grd devant $\operatorname{Card}\alpha$.]

% 19–20
Sous ces conditions, $I'$ est cofinal dans $I$ (donc \inserted{aussi} grand devant $\operatorname{Card}\alpha$, pour tt $\alpha\in I_{\alpha_0}$).

% 21–25
La démonstration est celle donnée \deleted{dans} ci-dessus. On n'a \unsure{pas eu} à utiliser que l'\deleted{\unsure{entier}}ordinal $\alpha_0=\alpha_{\unsure{\pi|\varpi}}$ \unsure{était} de cette espèce particulière, \unsure{seulement} que c'était un ordinal limite.
\scanpage{418}
\mspage{321}
% archive sheet number: 417

\deleted{

% 1
Il nous faut aussi une flèche de transition
% 2
\begingroup\csname tagsleft@true\endcsname
\[
\tau'_{\alpha,\alpha_0}:\unsure{j|J}_{\alpha_0}\longrightarrow j_\alpha,
\tag*{\unsure{(aaaa)|($\alpha\alpha\alpha\alpha$)|(****)}}
\]
\endgroup
% 3–9
mais comme $J'$ est filtrante, quitte à remplacer $j_\alpha$ par un majorant commun $j'_\alpha$ de $j_\alpha$ et de $j_{\alpha_0}$ dans $J$, on trouve \unsure{ainsi|aussi} (aaaa). Mais il faut voir si les quatre contraintes sont \unsure{vérifiées}, et au besoin modifier les choix précédents, de façon à ce qu'elles le soient.
}
\scanpage{419}
\mspage{331}
% archive sheet number: 418

% 1–5
Ce lemme implique, comme on l'a vu à la page 312, l'implication $b\Longrightarrow d$ dans le diagramme d'implications \unsure{(*)|($\alpha$)|(a)} page 310, lequel est donc à présent prouvé.

% 6–7
Pour achever de prouver le théorème 4.13.15, il reste à prouver les implications
% 8–9
\diagram{\begin{tikzcd}[ampersand replacement=\&,row sep=0.5em]
\deleted{\unsure{d}}\unsure{c'|d'}\arrow[dr,Rightarrow] \& \&\\
\& \deleted{\unsure{b}}\unsure{a'|b'}\arrow[r,Rightarrow] \& a)\\
b\arrow[ur,Rightarrow] \& \&
\end{tikzcd}}

% 10
Démontrons donc $[\unsure{b)|b'}\ \text{ou}\ \unsure{c')|d')}]\Longrightarrow a')$.

% 10–12, interline at right
\inserted{Il reste à prouver que $\underline S$ est $\unsure{\pi|\varpi}$-accessible dans $M$, i.e.} \deleted{Notons} \deleted{Comme $\underline S$ est déjà supposée str. pleine dans $M$, il nous reste à prouver} \deleted{que \unsure{$\underline S$}, \unsure{le} \illegible{1}}
% 12–15
que le foncteur pl. fidèle \deleted{$\pi$ \unsure{étant} \illegible{1} \unsure{pour} $M$,}
% 14
\begingroup\csname tagsleft@true\endcsname
\[
\operatorname{Ind}_\pi(\underline S(\pi))\longrightarrow\deleted{\unsure{\underline S}}\underline M
\tag*{(*)}
\]
\endgroup
% 15
a pour image essentielle \deleted{$\unsure{S}$} $\underline S$.

% 15–20
\deleted{$M_\pi/x$ est grande devant $\pi$, donc \inserted{la sous-catégorie str. pleine} $\underline S(\pi)/x$ est cofinale, elle est elle-$=$ grande devant $\pi$, a fortiori filtrante.

Donc c) implique que si $\unsure{x^-|x}\in\operatorname{Ob}\underline S$, alors on a $x=\varinjlim_{i\in\underline S(\pi)/x}x_i$, \deleted{ce qui implique} $\underline S(\pi)/x$ gde devant $\pi$. Comme}

% 21–22
Donc il faut prouver que si $x\in\operatorname{Ob} M$, alors
% 23–26
\begingroup\csname tagsleft@true\endcsname
\[
\begin{aligned}
 x\in\operatorname{Ob}\underline S\quad&\overset{\unsure{(?)|(!)}}{\Longleftrightarrow}\quad x\ \text{est dans l'image essentielle}\\
 &\hspace{4em}\text{de }\operatorname{Ind}_\pi(\underline S(\pi)),\ \text{i.e. }x\\
 &\hspace{4em}\text{est isomorphe à }\varinjlim_{i\in I}x_i,\\
 &\hspace{4em}\text{les }x_i\ \text{dans }\operatorname{Ob}(\underline S(\pi)\unsure{\,|)} ,\ I\ \text{grd}\\
 &\hspace{4em}\text{devant }\pi.
\end{aligned}
\tag*{$\substack{(*)\\\deleted{\unsure{(*)}}}$}
\]
\endgroup
\scanpage{420}
\mspage{332}
% archive sheet number: 419

% upper left, beside lines 1–4
% slanted
\margin{\unsure{Maladroit} \unsure{on peut|il faut} traiter les deux cas b) et \unsure{c')|a')} en $=$ temps cf. bas de \unsure{pge|pg} 333}

% top, lines 1–2
Supposons d'abord que $\underline S$ satisfasse b)\textsuperscript{*}. Donc on sait que
% upper third, lines 3–4
\begingroup\csname tagsleft@true\endcsname
\[
\begin{aligned}
x\in\operatorname{Ob}\underline S
  &\Longleftrightarrow \deleted{\unsure{\underline S}}\ \underline S(\pi)_{/x}\ \text{est \unsure{cofinal|cofinale}}\\
  &\hphantom{\Longleftrightarrow}\text{dans }M_{\pi}^{\deleted{\unsure{x}}}/x
\end{aligned}
\tag*{(**)}
\]
\endgroup

% upper third, lines 5–6
Nous voulons déduire l'équivalence (*) de l'équivalence (**).

% upper–middle, lines 7–12
Prouvons $\Rightarrow$ dans (*), i.e. supposons $x\in\operatorname{Ob}\underline S$, et prouvons que $x$ est dans l'image essentielle de $\operatorname{Ind}_{\pi}\underline S(\pi)$. En vertu de (**), $\underline S(\pi)_{/x}$ est \inserted{une sous-catégorie (str. pleine)} cofinale dans $M_{\pi}/x$, donc comme $M_{\pi}/x$ est \unsure{grd|grde} devant $\pi$, $\underline S(\pi)_{/x}$ l'est aussi. \deleted{\unsure{Donc enfin|Donc ceci}} Comme on sait que
% middle, lines 13–14
\[
x=\varinjlim_{i\in M_{\pi}/x}x_i
\]
% middle–lower, lines 15–16
($\pi$ étant cardinal utile pour $M$), il s'ensuit \deleted{que} par cofinalité qu'on a aussi
% lower third, lines 17–18
\[
x=\varinjlim_{i\in\underline S(\pi)_{/x}}x_{\unsure{.|i}}, 
\]
% lower third, lines 19–21
et comme $\underline S(\pi)_{/x}$ est grd devant $\pi$, cela implique que $x$ appartient à l'image essentielle de $\operatorname{Ind}_{\pi}(\underline S(\pi))$.

% foot, last two lines
Prouvons $\Leftarrow$ dans (*), i.e. supposons que $x$ soit dans l'image essentielle de $\operatorname{Ind}_{\pi}(\underline S(\pi))$,
